\documentclass[11pt,a4paper]{article}

% Packages
\usepackage[UTF8]{ctex}
\usepackage{amsmath,amssymb,amsfonts}
\usepackage{graphicx}
\usepackage{hyperref}
\usepackage{listings}
\usepackage{xcolor}
\usepackage{booktabs}
\usepackage{geometry}
\usepackage{fancyhdr}
\usepackage{longtable}
\usepackage{array}
\usepackage{enumitem}
\usepackage{caption}
\usepackage{subcaption}
\usepackage{multirow}

\geometry{margin=2.5cm}
\hypersetup{
    colorlinks=true,
    linkcolor=blue,
    citecolor=blue,
    urlcolor=blue,
}

% Code listing style
\lstdefinestyle{pythonstyle}{
    language=Python,
    basicstyle=\ttfamily\footnotesize,
    keywordstyle=\color{blue}\bfseries,
    commentstyle=\color{gray}\itshape,
    stringstyle=\color{red},
    numbers=left,
    numberstyle=\tiny\color{gray},
    frame=single,
    breaklines=true,
    tabsize=2,
    showstringspaces=false,
}
\lstdefinestyle{bashstyle}{
    language=bash,
    basicstyle=\ttfamily\footnotesize,
    frame=single,
    breaklines=true,
    showstringspaces=false,
}

% Severity colors
\definecolor{severityfatal}{HTML}{CC0000}
\definecolor{severityhigh}{HTML}{DD4400}
\definecolor{severitymed}{HTML}{CC8800}
\definecolor{severitylow}{HTML}{228822}
\definecolor{severityperf}{HTML}{2266CC}
\newcommand{\fatal}{\textcolor{severityfatal}{\textbf{FATAL}}}
\newcommand{\high}{\textcolor{severityhigh}{\textbf{HIGH}}}
\newcommand{\med}{\textcolor{severitymed}{\textbf{MEDIUM}}}
\newcommand{\low}{\textcolor{severitylow}{\textbf{LOW}}}
\newcommand{\perf}{\textcolor{severityperf}{\textbf{PERF}}}

\title{\textbf{PyQCU 开发报告}\\
\large stab22 $\rightarrow$ dev73}
\author{Zhang Xin}
\date{2026-07-28}

\begin{document}

\maketitle
\tableofcontents
\newpage

%==========================================================================
\section{时间线概览}
%==========================================================================

本报告对照 tag \texttt{stab22}（2026-07-05），总结截至 2026-07-28 的 PyQCU 全库变更，
涵盖性能优化、新功能（C++ CUDA Clover 多网格求解器）、Bug 修复、三轮全面代码审查、文档更新及测试验证。

总代码变更：\textbf{68 files, +9,282 / $-$333 lines}。

\begin{table}[ht]
\centering
\caption{开发时间线}
\begin{tabular}{lll}
\toprule
日期 & 里程碑 & 关键变更 \\
\midrule
07-05 & \textbf{stab22} & 基线标签 \\
07-08 & 性能优化 & 12 项优化，6 文件，+515/$-$110 行 \\
07-08 & CLAUDE.md 初始化 & 90 行项目文档 \\
07-14 & CLAUDE.md 扩展 & +36 行，补充架构细节 \\
07-17 & CLQCD 报告 & \texttt{refer/dev71.md}(861行) + .pdf + .tex(853行) \\
07-17 & C++ CUDA MG 骨架 & 多网格类 + Python CUDA 桥接 \\
07-19 & MG 后端修复 & 重构 MPI 通信 + 流同步修复 \\
07-20 & 基础设置 & PASS 标记, past-work.md, CLAUDE.md 大修 \\
\textbf{07-28} & \textbf{大规模审查+修复} & Clover MG 求解器, 3 轮审查 135 项发现, 33 Bug 修复 \\
\bottomrule
\end{tabular}
\end{table}

%==========================================================================
\section{性能优化（07-08）}
%==========================================================================

\textbf{来源：} \texttt{OPTIMIZATION\_REPORT.md} $\rightarrow$ \texttt{log/stab23.log}（commit \texttt{a3c3e96}）\\
\textbf{修改文件：} 6 个 \quad \textbf{代码变更：} +120 / $-$110 行

\subsection{批量矩阵求逆}

\texttt{pyqcu/dslash/\_clover.py} 中 \texttt{inverse()} 原使用 Python for 循环逐一求 $12\times12$ 复矩阵的逆。
优化后调用 \texttt{torch.linalg.inv} 批量处理，将 $N$ 次独立 kernel launch 合并为 1 次，预计 L=32 格点上 \textbf{10--50x} 加速。

\begin{lstlisting}[style=pythonstyle]
# 优化前：逐格点循环
for i in range(_clover_term.shape[-1]):
    _clover_term[:, :, i] = torch.linalg.inv(_clover_term[:, :, i])

# 优化后：批量求逆
_clover_term = torch.linalg.inv(
    _clover_term.permute(2, 0, 1)).permute(1, 2, 0)
\end{lstlisting}

\subsection{张量设备/类型缓存}

\texttt{pyqcu/dslash/\_wilson.py} 中在 Wilson dslash 入口处预计算 $I \pm \gamma$ 矩阵字典，
消除 4 方向 $\times$ 多次 \texttt{.to()/.type()} 的 GPU 显存分配。同样应用于
\texttt{give\_wilson\_eo}, \texttt{give\_wilson\_oe}, \texttt{give\_hopping\_plus}, \texttt{give\_hopping\_minus}。

\subsection{预计算 sigma 矩阵与 clover 系数}

\begin{itemize}
    \item \textbf{sigma 矩阵缓存：} 6 个 $\gamma\gamma$ 矩阵预计算到设备，消除 6 次重复 \texttt{.to()} 调用
    \item \textbf{clover 系数：} \texttt{float(kappa/u\_0)} 从循环内提到循环外
\end{itemize}

\subsection{内存优化}

移除 \texttt{add\_I()}, \texttt{cut\_I()}, \texttt{inverse()} 中不必要的 \texttt{.clone()} 调用。
每次调用节省约 72\,MB（$12\times12\times32^4\times8$ bytes）内存分配。

\subsection{缓存 \texttt{give\_eo\_mask}}

引入全局缓存字典，按 \texttt{(Lx, Ly, Lz, Lt, eo, device)} 键缓存 checkerboard mask，避免重复 \texttt{meshgrid} 创建。

\subsection{求解器优化}

\begin{itemize}
    \item 存储 \texttt{norm(b)} 避免重复 MPI Allreduce
    \item \texttt{verbose=False} 时跳过 \texttt{perf\_counter()} 调用
    \item 移除 \texttt{\_multigrid.py} 重复 import
    \item 移除 \texttt{\_linalg.py} 中 \texttt{norm()} 的冗余 \texttt{.flatten()}
\end{itemize}

\subsection{预期性能提升}

\begin{table}[ht]
\centering
\caption{性能优化预期效果}
\begin{tabular}{ll}
\toprule
场景 & 预期加速 \\
\midrule
Clover 求逆 (L=32) & \textbf{10--50x} \\
Wilson Dslash 单次 & \textbf{1.5--2x} \\
Clover 项构建 & \textbf{1.2--1.5x} \\
BiStabCG silent 模式 & 每次迭代 $-2\%$ \\
Parity Dslash & mask 计算 $-5\%$ \\
Clover add\_I/cut\_I & 每次调用 $-70$\,MB 峰值 \\
\bottomrule
\end{tabular}
\end{table}

%==========================================================================
\section{CLQCD 报告与参考资料（07-17）}
%==========================================================================

\textbf{来源：} \texttt{refer/dev71.md}, \texttt{refer/dev71.tex}, \texttt{refer/dev71.pdf}（commits \texttt{fcf6e9d}, \texttt{34c71c6}）

创建 CLQCD 会议/论文参考资料，标题\textbf{``CUDA C++ 版 MultiGrid 编写与优化解析文档''}
（861 行 Markdown + 853 行 \LaTeX{} + 编译 PDF 345\,KB）。

内容涵盖：
\begin{itemize}
    \item 总体架构概览（Python/C++ 双层设计）
    \item Python 层 MultiGrid 完整实现分析（336 行源码步行）
    \item C++ CUDA 后端现有基础设施（\texttt{lattice\_set.h}, \texttt{lattice\_wilson\_bistabcg.h}, \texttt{lattice\_clover\_bistabcg.h} 等）
    \item CUDA C++ MultiGrid 设计方案（V-cycle 流程、流并行策略、通信模式）
    \item 优化策略（内存合并、共享内存、warp shuffle、寄存器优化）
    \item 实现路线图（Phase 1--4）
\end{itemize}

该文档作为后续 C++ Clover Multigrid 实现的技术规格。

%==========================================================================
\section{C++ CUDA 多网格后端（07-17 $\rightarrow$ 07-28）}
%==========================================================================

\subsection{初始骨架（07-17）}

\textbf{新增文件：}
\begin{itemize}
    \item \texttt{cpp/cuda/qcu/include/lattice\_multigrid.h} — 59 行，多网格类声明
    \item \texttt{cpp/cuda/qcu/include/multigrid.h} — 20 行，multigrid 参数结构体
    \item \texttt{cpp/cuda/qcu/src/multigrid.cu} — 232 行，V-cycle + 粗网格算子实现
    \item \texttt{cpp/cuda/qcu/src/apply\_multigrid.cu} — 125 行，C API 桥接函数
\end{itemize}

\textbf{核心能力：}
\begin{itemize}
    \item \texttt{applyMultigridRestrictQcu} — CUDA 限制算子
    \item \texttt{applyMultigridProLongQcu} — CUDA 延拓算子
    \item \texttt{applyMultigridCoarseDslashQcu} — CUDA 粗网格 Dirac 算子
    \item Python 层 \texttt{with\_cuda\_qcu=True} 自动启用 CUDA 加速
\end{itemize}

\subsection{Clover Multigrid 求解器（07-28）}

\textbf{重大新增：} 完整的 C++ CUDA Clover 多网格求解器。

新增核心文件：
\begin{itemize}
    \item \texttt{cpp/cuda/qcu/include/lattice\_clover\_multigrid.h} — 953 行，主求解器类
    \item \texttt{cpp/cuda/qcu/src/apply\_clover\_multigrid.cu} — 72 行，C API 桥接
    \item \texttt{examples/qcu/conftest.clover.multigrid.py} — 测试驱动脚本
\end{itemize}

\subsubsection{关键设计决策}

\begin{table}[ht]
\centering
\caption{Clover Multigrid 求解器架构}
\begin{tabular}{ll}
\toprule
特性 & 说明 \\
\midrule
流并行 & 5 个 CUDA 流（\texttt{strm} + \texttt{\_a\_}/\texttt{\_b\_}/\texttt{\_c\_}/\texttt{\_d\_}） \\
标量管理 & \texttt{device\_vals} 仅被 GPU kernel 修改，无 host$\rightarrow$device memcpy \\
同步策略 & 每轮迭代底部 5 流全同步，顶部 5 流全同步 \\
dot 协议 & cublasDot $\rightarrow$ \texttt{\_send\_tmp\_} $\rightarrow$ MPI\_Allreduce $\rightarrow$ 目标槽位 \\
精度模板 & \texttt{mpi\_real\_type<T>()} 根据 float/double 选择 MPI\_FLOAT/MPI\_DOUBLE \\
\bottomrule
\end{tabular}
\end{table}

\subsubsection{测试结果}

\textbf{GPU：} NVIDIA GeForce RTX 4060 Laptop (SM 8.9)

\begin{table}[ht]
\centering
\caption{Clover Multigrid 正确性与性能}
\begin{tabular}{ll}
\toprule
指标 & 数值 \\
\midrule
正确性 ($\|x_{\text{mg}}-x_{\text{ref}}\|/\|x_{\text{ref}}\|$) & \textbf{$6.23\times10^{-7}$} \\
NaN 出现次数 & \textbf{0} \\
收敛迭代次数 & $\sim$104 到 $5.7\times10^{-7}$ \\
vs BiStabCG & 0.87$\times$--1.28$\times$（随运行变化）\\
求解时间 & 1742--3636 ms（冷启动$\rightarrow$稳定）\\
\bottomrule
\end{tabular}
\end{table}

\subsubsection{Clover MG Bug 修复}

\begin{table}[ht]
\centering
\caption{Clover Multigrid 求解器 Bug 修复清单}
\begin{tabular}{cllp{4cm}}
\toprule
\# & 严重度 & Bug & 修复 \\
\midrule
1 & \fatal & \texttt{run\_mpi} 未初始化 MPI\_Wait & 移除阻塞路径的 MPI\_Wait \\
2 & \fatal & \texttt{device\_vals} 竞争条件 $\rightarrow$ NaN & 移除迭代内 host$\rightarrow$device memcpy \\
3 & \high & 底部流同步缺失 & 底部 5 流全同步 \\
4 & \med & cublasDot 目标槽位损坏 & 先写 \texttt{\_send\_tmp\_}，MPI 后复制 \\
5 & \high & MPI\_FLOAT 硬编码 & \texttt{mpi\_real\_type<T>()} 模板 \\
6 & \fatal & \texttt{wilson\_dslash.cu} 变量名错误 & \texttt{idx} $\rightarrow$ \texttt{parity} (8处) \\
\bottomrule
\end{tabular}
\end{table}

%==========================================================================
\section{基础设施与文档（07-14 $\rightarrow$ 07-28）}
%==========================================================================

\subsection{CLAUDE.md 演进}

\begin{table}[ht]
\centering
\caption{CLAUDE.md 迭代历史}
\begin{tabular}{lll}
\toprule
日期 & 行数 & 主要变化 \\
\midrule
07-08 & 90 & 初始化：项目结构、构建命令、架构概述 \\
07-14 & +36 & GPU 后端抽象、测试运行方式、模块级代码说明 \\
07-20 & +183 (大修) & 两层架构图、数据布局约定、C++ 计划系统、Cython 桥接、TileLang、cann 层 \\
07-28 & +66 & 方法列表细化、cann 详细文档、GMRES stub 状态、ward 索引惯例 \\
\bottomrule
\end{tabular}
\end{table}

\subsection{其他基础设施}

\begin{itemize}
    \item \textbf{\texttt{.gitignore}} (07-17): 新增 9 条忽略规则
    \item \textbf{\texttt{cpp/*/PASS}} (07-20): CANN/CUDA/DTK/MACA 四个后端放置 PASS 标记文件
    \item \textbf{\texttt{.claude/skills/past-work.md}} (07-20): 115 行历史工作总结
    \item \textbf{\texttt{build.sh} / \texttt{make.sh}} (07-28): 添加 \texttt{set -e} 错误检测、\texttt{\&\&} 链式调用、\texttt{rm -f} 安全删除
    \item \textbf{\texttt{pyqcu/cuda/\_\_init\_\_.py}} (07-28, 新建): 修复 pip install 后 CUDA 包不可用问题
    \item \textbf{\texttt{pyqcu/cuda/qcu/qcu.pyi}} (07-28, 新建): 155 行类型桩文件，包含所有 22 个 Cython 桥接函数签名
\end{itemize}

%==========================================================================
\section{全面代码审查（07-28）}
%==========================================================================

\subsection{审查流程}

\begin{table}[ht]
\centering
\caption{三轮审查流程}
\begin{tabular}{ll}
\toprule
轮次 & 流程 \\
\midrule
Round 1 (74 findings) & $\rightarrow$ 12 bug fixes + 1 optimization \\
Round 2 (28 findings) & $\rightarrow$ 8 bug fixes + 1 optimization + 5 refinements + 6 docs \\
Round 3 (33 findings) & $\rightarrow$ 13 fixes + 2 docs \\
\midrule
\textbf{合计 135 findings} & \textbf{$\rightarrow$ 33 bugs fixed + 2 optimizations + 12 docs + 13 false positives} \\
\bottomrule
\end{tabular}
\end{table}

\subsection{审查方法}

\begin{itemize}
    \item \textbf{Round 1：} 逐文件阅读，交叉引用 CLAUDE.md 文档
    \item \textbf{Round 2：} 三轮并行 agent 审查 + 人工深度追踪关键算法路径（Python core + C++ backend + Cython/tests/docs + deep-dive）
    \item \textbf{Round 3：} 四路并行 agent（跨模块一致性 + 物理/数值正确性 + 构建/测试/边界 + 人工 multigrid cycle 追踪）
\end{itemize}

\subsection{审查报告文件}

\begin{table}[ht]
\centering
\caption{审查与修复文档}
\begin{tabular}{lll}
\toprule
文件 & 行数 & 描述 \\
\midrule
\texttt{review-2026-07-28.md} & 892 & Round 1 审查报告 \\
\texttt{review-2026-07-28-r2.md} & 460 & Round 2 审查报告 \\
\texttt{review-2026-07-28-r3.md} & 460 & Round 3 审查报告 \\
\texttt{fix-report-2026-07-28.md} & 210 & 完整修复清单 \\
\texttt{debug/fix-log.md} & 48 & R1 修复日志 \\
\texttt{debug/fix-log-r2.md} & 82 & R2 修复日志 \\
\texttt{debug/fix-log-r3.md} & 123 & R3 修复日志 \\
\bottomrule
\end{tabular}
\end{table}

\subsection{发现的假阳性（误报）}

13 项经深度分析确认为设计决策而非 Bug：

\begin{table}[ht]
\centering
\caption{假阳性确认清单}
\begin{tabular}{p{4cm}p{5cm}p{4.5cm}}
\toprule
项目 & 审查发现 & 分析结论 \\
\midrule
Gamma 矩阵代数 & 怀疑平方/反对易关系错误 & \checkmark 所有 5 个 gamma 恒等式正确 \\
Wilson dslash 符号 & 怀疑方向索引排列错误 & \checkmark 所有 22 个 einsum 方程正确 \\
BiCGStab 算法 & 怀疑与 van der Vorst 1992 不符 & \checkmark 精确匹配标准算法 \\
MG V-cycle & 怀疑 Galerkin 修正错误 & \checkmark 修正和分解逻辑完全正确 \\
C++ host\_vals 同步 & 怀疑竞争条件 & \checkmark D2H$\rightarrow$MPI$\rightarrow$H2D 顺序正确 \\
BiCGStab 缓冲区泄漏 & 怀疑 GPU buffer 泄漏 & \checkmark 在 \texttt{end()} 中正确释放 \\
MPI\_FLOAT 硬编码 & 怀疑精度错误 & \checkmark 已在早期修复 \\
导入图 & 怀疑循环导入 & \checkmark 无环（部分模块加载）\\
.pxd vs .pyx & 怀疑签名不匹配 & \checkmark 所有 22 个函数签名一致 \\
ward 负索引 & 怀疑索引错误 & \checkmark 有意的设计模式 \\
\bottomrule
\end{tabular}
\end{table}

%==========================================================================
\section{Bug 修复总表}
%==========================================================================

\subsection{\fatal{} 致命 Bug（5 个，全部修复）}

\begin{table}[ht]
\centering
\caption{致命 Bug 修复清单}
\begin{tabular}{clp{5cm}p{4cm}l}
\toprule
\# & 文件 & 问题 & 症状 & 轮次 \\
\midrule
F1 & \texttt{lattice\_complex.h:79-83} & \texttt{operator*=} 复数乘法覆盖 bug & 所有复数乘法结果错误 & R1 \\
F2 & \texttt{gauss\_gauge.cu:183-186} & 非 verbose 路径内存分配不足 & 越界写入 (OOB Write) & R1 \\
F3 & \texttt{lattice\_wilson\_cg.h:41-51} & CG \texttt{\_init()} 重复分配缓冲区 & GPU 内存泄漏 & R2 \\
F4 & \texttt{lattice\_set.h:140-163} & Grid dim 整数除法截断 & 格点遗漏 & R2 \\
F5 & \texttt{pyqcu/cuda/} 无 \texttt{\_\_init\_\_.py} & pip install 后包不可用 & ImportError & R3 \\
\bottomrule
\end{tabular}
\end{table}

\subsection{\high{} 严重 Bug（14 个，全部修复）}

\begin{longtable}{cll}
\toprule
\# & 文件 & 问题 \\
\midrule
\endfirsthead
\multicolumn{3}{c}{\tablename{} \thetable{} —— 续} \\
\midrule
\# & 文件 & 问题 \\
\midrule
\endhead
\bottomrule
\endfoot
S1 & \texttt{\_io.py:62,69} & gather 元组 (t,z,y,x) $\rightarrow$ unpack (x,y,z,t) 索引错位 \hfill R1 \\
S2 & \texttt{\_stout.py:64-172} & \texttt{nstep>1} 无效：每次迭代使用原始 U \hfill R1 \\
S3 & \texttt{\_stout.py:22-63} & MPI 边界数据在多次 smearing 中不更新 \hfill R1 \\
S4 & \texttt{cann/\_\_init\_\_.py:81-83} & NPU 3+ operand 复数 einsum 丢弃虚部 \hfill R1 \\
S5 & \texttt{\_operator.py:228,259} & \texttt{self.sitting} 对象始终 truthy \hfill R1 \\
S6 & \texttt{\_define.py:97-101} & \texttt{check\_mpi\_support} 非 root 进程临时文件泄漏 \hfill R1 \\
S7 & \texttt{testing/\_\_init\_\_.py:363} & 错误消息使用模块对象而非参数 \hfill R1 \\
S8 & \texttt{setup.py:1,3} & distutils 覆盖 setuptools \texttt{setup} \hfill R1 \\
S9 & \texttt{qcu.pyx:2-13} & 模块级 cdef 变量 GIL 依赖风险 \hfill R1 \\
S10 & \texttt{\_operator.py:333-334} & 遗留调试 print() \hfill R1 \\
S11 & \texttt{\_stout.py:131-136} & $c_1=0 \rightarrow c_{0,\max}=0 \rightarrow \arccos(0/0) \rightarrow$ NaN \hfill R2 \\
S12 & \texttt{\_stout.py:149} & f\_denom 除以零（$9u^2=w^2$ 时）\hfill R2 \\
S13 & \texttt{\_stout.py:155-162} & NPU f1 宇称缺失实部取反 \hfill R2 \\
S14 & \texttt{smear/\_stout.py:160-172} & R2 f1/f2 宇称修复虚部符号反转回归 \hfill R3 \\
\end{longtable}

\subsection{\med{} 中等 Bug（13 个，全部修复）}

\begin{longtable}{cll}
\toprule
\# & 文件 & 问题 \\
\midrule
\endfirsthead
\multicolumn{3}{c}{\tablename{} \thetable{} —— 续} \\
\midrule
\# & 文件 & 问题 \\
\midrule
\endhead
\bottomrule
\endfoot
M1 & \texttt{\_bistabcg.py:66} & \texttt{verbose=False} 时 ZeroDivisionError \hfill R1 \\
M2 & \texttt{\_bistabcg.py:66-72} & 性能统计无视 verbose 标志 \hfill R1 \\
M3 & \texttt{cuda/define.py:94-119} & \texttt{dtype()} 裸 raise 无异常消息 \hfill R1 \\
M4 & \texttt{setup.py:38} & exclude pattern 不匹配 \texttt{pyqcu.testing} \hfill R1 \\
M5 & \texttt{apply\_end.cu} & LatticeSet 对象未 delete \hfill R1 \\
M6 & \texttt{lattice\_wilson\_dslash.h} & MPI\_Isend 从未 MPI\_Wait \hfill R1 \\
M7 & \texttt{lattice\_complex.h:89-95} & \texttt{operator/=} 与 \texttt{operator*=} 同类 bug \hfill R1 \\
M8 & \texttt{gauss\_gauge.cu} & GPU 内存泄漏（无对应 cudaFreeAsync）\hfill R1 \\
M9 & \texttt{\_linalg.py:21,26} & 冗余 MPI.Barrier 包围 Allreduce \hfill R2 \\
M10 & \texttt{\_bistabcg.py:38-47} & BiCGStab 无除以零检测 \hfill R2 \\
M11 & \texttt{\_multigrid.py:351-359} & MG BiCGStab 同上除以零问题 \hfill R2 \\
M12 & \texttt{\_multigrid.py:407-420} & MG cycle 粗网格修正后状态未重置 \hfill R3 \\
M13 & \texttt{setup.py:46} & \texttt{python\_requires ">=3.6"} 与 PyTorch 2.x 不兼容 \hfill R3 \\
\end{longtable}

\subsection{\low{} 代码质量（6 项修复）}

\begin{longtable}{cl}
\toprule
\# & 问题 \\
\midrule
\endfirsthead
\multicolumn{2}{c}{\tablename{} \thetable{} —— 续} \\
\midrule
\# & 问题 \\
\midrule
\endhead
\bottomrule
\endfoot
Q1 & \texttt{lattice/\_\_init\_\_.py} 重复 \texttt{from typing import Optional} \hfill R1 \\
Q2 & 多处使用可变默认参数 \texttt{torch.Tensor([0.1])} \hfill R1 \\
Q3 & Gamma 矩阵 ward 负索引惯例添加注释 \hfill R1 \\
Q4 & \texttt{check\_su3} float32 默认 \texttt{atol=1e-8} 过严 \hfill R2 \\
Q5 & \texttt{build.sh} / \texttt{make.sh} 无 \texttt{set -e} 错误检测 \hfill R3 \\
Q6 & 多处 bare except $\rightarrow$ except Exception \hfill R3 \\
\end{longtable}

\subsection{\perf{} 性能优化（2 项）}

\begin{longtable}{cl}
\toprule
\# & 优化 \\
\midrule
\endfirsthead
\bottomrule
\endfoot
P1 & 移除 3 个文件中 24 个冗余 MPI.Barrier()（阻塞 Sendrecv 不需要 Barrier）\hfill R1 \\
P2 & ortho\_r/ortho\_null\_vecs 当 normalize=True 跳过冗余 vdot 分母 \hfill R2 \\
\end{longtable}

\subsection{测试与类型桩修复（3 项）}

\begin{longtable}{cl}
\toprule
\# & 修复 \\
\midrule
\endfirsthead
\bottomrule
\endfoot
T1 & \texttt{pyqcu/cuda/qcu/qcu.pyi} — 155 行类型桩，覆盖 22 个函数 \hfill R3 \\
T2 & \texttt{pyqcu/testing/\_\_init\_\_.py} — 添加 pytest assert（原只有 print）\hfill R3 \\
T3 & \texttt{examples/profiler/conftest.py} — import comm $\rightarrow$ from mpi4py import MPI \hfill R3 \\
\end{longtable}

%==========================================================================
\section{文件变更统计}
%==========================================================================

\subsection{总体统计}

\begin{verbatim}
68 files changed, 9282 insertions(+), 333 deletions(-)
\end{verbatim}

\subsection{新增文件（20+ 个）}

\begin{longtable}{llr}
\toprule
类别 & 文件 & 行数 \\
\midrule
\endfirsthead
\multicolumn{3}{c}{\tablename{} \thetable{} —— 续} \\
\midrule
类别 & 文件 & 行数 \\
\midrule
\endhead
\bottomrule
\endfoot
\multirow{7}{*}{C++ 后端} & \texttt{lattice\_clover\_multigrid.h} & 953 \\
& \texttt{lattice\_multigrid.h} & 59 \\
& \texttt{multigrid.h} & 20 \\
& \texttt{apply\_clover\_multigrid.cu} & 72 \\
& \texttt{apply\_multigrid.cu} & 125 \\
& \texttt{multigrid.cu} & 232 \\
& \texttt{wilson\_dslash.cu.bak} & 1462 \\
\midrule
\multirow{19}{*}{文档} & \texttt{log/bug30.md} & 892 \\
& \texttt{log/review-2026-07-28.md} & 838 \\
& \texttt{log/review-2026-07-28-r2.md} & 460 \\
& \texttt{log/review-2026-07-28-r3.md} & 460 \\
& \texttt{log/fix-report-2026-07-28.md} & 210 \\
& \texttt{log/multigrid\_report.md} & 144 \\
& \texttt{log/debug/fix-log.md} & 48 \\
& \texttt{log/debug/fix-log-r2.md} & 82 \\
& \texttt{log/debug/fix-log-r3.md} & 123 \\
& \texttt{log/results/final-report.md} & 55 \\
& \texttt{log/results/remaining-fixes-report.md} & 45 \\
& \texttt{refer/dev71.md} & 861 \\
& \texttt{refer/dev71.tex} & 853 \\
& \texttt{CLAUDE.md} (扩充) & 271 \\
& \texttt{log/dev73.md} (本报告) &  \\
& \texttt{log/multigrid\_performance.png} & 59\,KB \\
& \texttt{log/multigrid\_performance\_0.png} & 218\,KB \\
& \texttt{log/multigrid\_result.png} & 73\,KB \\
& \texttt{refer/dev71.pdf} & 345\,KB \\
\midrule
\multirow{2}{*}{Cython/类型} & \texttt{pyqcu/cuda/qcu/qcu.pyi} & 155 \\
& \texttt{pyqcu/cuda/\_\_init\_\_.py} & 10 \\
\midrule
测试 & \texttt{examples/qcu/conftest.clover.multigrid.py} & 138 \\
\midrule
\multirow{2}{*}{配置} & \texttt{.claude/skills/past-work.md} & 115 \\
& \texttt{.gitignore} & 9 \\
\end{longtable}

\subsection{修改核心文件}

\begin{longtable}{lrl}
\toprule
文件 & 变更 & 描述 \\
\midrule
\endfirsthead
\multicolumn{3}{c}{\tablename{} \thetable{} —— 续} \\
\midrule
文件 & 变更 & 描述 \\
\midrule
\endhead
\bottomrule
\endfoot
\texttt{pyqcu/dslash/\_wilson.py} & +170 & 性能优化 + 缓存 \\
\texttt{pyqcu/smear/\_stout.py} & +144 & nstep 修复 + NaN 防护 + MPI 边界 \\
\texttt{pyqcu/solver/\_multigrid.py} & +187 & CUDA 加速 + 除以零检测 + 状态重置 \\
\texttt{pyqcu/dslash/\_operator.py} & +63 & 条件修复 + 调试 print 删除 \\
\texttt{pyqcu/dslash/\_clover.py} & +58 & 批量求逆 + 预计算 \\
\texttt{pyqcu/cann/\_\_init\_\_.py} & +55 & 3+ operand einsum \\
\texttt{pyqcu/lattice/\_\_init\_\_.py} & +53 & check\_su3 修复 + ward 注释 \\
\texttt{pyqcu/cuda/qcu/qcu.pyx} & +51 & MG 桥接函数 + cdef 声明 \\
\texttt{pyqcu/solver/\_bistabcg.py} & +45 & norm(b) 缓存 + 除以零检测 \\
\texttt{pyqcu/tools/\_multigrid.py} & +42 & 正交化优化 \\
\texttt{pyqcu/tools/\_define.py} & +31 & eo\_mask 缓存 + 异常处理 \\
\texttt{pyqcu/tools/\_io.py} & +24 & 索引修复 \\
\texttt{pyqcu/testing/\_\_init\_\_.py} & +20 & assert + 异常处理 \\
\texttt{pyqcu/cuda/define.py} & +17 & dtype ValueError \\
\texttt{setup.py} & +13 & python\_requires + exclude \\
\end{longtable}

%==========================================================================
\section{测试验证结果}
%==========================================================================

\subsection{Round 1 验证（8/8 PASSED）}

\begin{verbatim}
============================================================
  PyQCU Validation Suite — 2026-07-28
============================================================
  PASS: stout_smear nstep>1
  PASS: operator parity
  PASS: BiCGStab
  PASS: BiCGStab parity
  PASS: NPU 2-op einsum
  PASS: NPU 3-op einsum
  PASS: No MPI orphans
  PASS: cuda/define ValueError

  TOTAL: 8/8 passed — ALL TESTS PASSED
\end{verbatim}

C++ 构建：\texttt{[100\%] Linking CUDA shared library libqcu.so} \checkmark

\subsection{Round 2 验证（8/8 + 10/10 PASSED）}

\begin{verbatim}
PASS: check_su3 float32 (atol fix)
PASS: stout NaN guard (trivial gauge)
PASS: stout_smear nstep>1
PASS: operator parity
PASS: BiCGStab + breakdown guard
PASS: NPU stout parity fix
PASS: NPU einsum 3-op
PASS: vdot after Barrier removal

Extra: PASS: set_device verbose=False
       PASS: NPU restrict validation
       PASS: ortho_r vdot cache (safe matvec)
\end{verbatim}

\subsection{Round 3 验证（13/13 PASSED）}

\begin{verbatim}
PASS: R3 NPU stout f1/f2 fix
PASS: pyqcu.cuda __init__.py
PASS: check_su3 atol
PASS: stout NaN guard
PASS: stout nstep>1
PASS: operator parity
PASS: BiCGStab
PASS: MG solve
PASS: test_lattice with assert
PASS: dtype() raises ValueError
PASS: No bare except in _define.py
PASS: vdot Barrier removal
PASS: set_device verbose=False

13/13 passed — ALL R3 FIXES VERIFIED
\end{verbatim}

C++ 构建：\texttt{libqcu.so: 22.8 MB} \checkmark

\subsection{Clover Multigrid 正确性验证}

\begin{table}[ht]
\centering
\caption{Clover MG 最终验证}
\begin{tabular}{ll}
\toprule
指标 & 结果 \\
\midrule
相对残差 ($\|x_{\text{mg}}-x_{\text{ref}}\|/\|x_{\text{ref}}\|$) & $6.23\times10^{-7}$ \checkmark \\
NaN 出现次数 & 0 \checkmark \\
收敛到 $5.7\times10^{-7}$ 迭代数 & $\sim$104 \\
日志格式 & 匹配 v20260506 参考日志 \checkmark \\
\bottomrule
\end{tabular}
\end{table}

%==========================================================================
\section{参考图与参考源}
%==========================================================================

\subsection{图表}

\begin{table}[ht]
\centering
\caption{性能与收敛图表}
\begin{tabular}{lll}
\toprule
文件 & 描述 & 大小 \\
\midrule
\texttt{multigrid\_performance.png} & Clover MG 性能对比 + 收敛历史 & 59\,KB \\
\texttt{multigrid\_performance\_0.png} & 初始性能测试（高分辨率）& 218\,KB \\
\texttt{multigrid\_result.png} & 最终性能对比 + 收敛历史 & 73\,KB \\
\bottomrule
\end{tabular}
\end{table}

\subsection{关键源码}

\begin{itemize}
    \item \texttt{cpp/cuda/qcu/include/lattice\_clover\_multigrid.h} — Clover MG 求解器源码（最终版本 1128 行）
    \item \texttt{examples/qcu/conftest.clover.multigrid.py} — MG 测试驱动脚本
    \item \texttt{pyqcu/cuda/qcu/qcu.pyi} — 155 行类型桩（22 函数）
    \item \texttt{pyqcu/cuda/\_\_init\_\_.py} — 包初始化修复
\end{itemize}

\subsection{原始日志}

\begin{itemize}
    \item \texttt{clover\_multigrid.log} — C++ 收敛日志（\texttt{PYQCU::SOLVER::MULTIGRID::} 格式）
    \item \texttt{clover\_multigrid\_report.log} — 性能报告日志（13 次测试运行）
    \item \texttt{multigrid\_report.json} — JSON 格式收敛数据
    \item \texttt{test\_multigrid.log} — Python 测试执行日志
\end{itemize}

\subsection{参考文档}

\begin{itemize}
    \item \texttt{refer/dev71.md} — CUDA C++ MultiGrid 编写与优化解析（861 行）
    \item \texttt{refer/dev71.tex} — \LaTeX{} 源文件（853 行）
    \item \texttt{refer/dev71.pdf} — 编译 PDF（345\,KB）
    \item \texttt{CLAUDE.md} — 项目 AI 助手文档（271 行）
\end{itemize}

%==========================================================================
\section{三轮修复总结}
%==========================================================================

\begin{table}[ht]
\centering
\caption{三轮审查与修复总计}
\begin{tabular}{lrrrrrr}
\toprule
轮次 & 审查发现 & Bug 修复 & 优化 & 文档 & 假阳性 & 累计修复 \\
\midrule
R1 & 74 & 12 & 1 & 4 & 0 & 13 \\
R2 & 28 & 8 & 1 & 4 & 3 & 22 \\
R2 剩余 & — & 0 & 5 & 6 & — & 27 \\
R3 & 33 & 13 & 0 & 2 & 10 & 40 \\
\midrule
\textbf{合计} & \textbf{135} & \textbf{33} & \textbf{7} & \textbf{16} & \textbf{13} & \textbf{40} \\
\bottomrule
\end{tabular}
\end{table}

\vspace{1cm}
\begin{center}
\rule{0.5\textwidth}{0.4pt}\\[4pt]
\textit{报告生成: 2026-07-28 \quad | \quad 从 stab22 (2026-07-05) 到 dev73 (2026-07-28)}\\[4pt]
\textit{总代码变更: 68 files, +9,282 / $-$333 lines}\\[4pt]
\textit{标签: stab23}
\end{center}

\end{document}
